% Guia de practica: Pokedex con Ionic 9 + Angular 22. Compilar con pdfLaTeX.
\documentclass[11pt,a4paper]{article}
\usepackage[utf8]{inputenc}
\usepackage[T1]{fontenc}
\usepackage{lmodern}
\usepackage[spanish,es-tabla]{babel}
\usepackage[margin=2.3cm]{geometry}
\usepackage{microtype}
\usepackage{amssymb}
\usepackage{array}
\usepackage{booktabs}
\usepackage{tabularx}
\usepackage{xcolor}
\usepackage{enumitem}
\usepackage{fancyhdr}
\usepackage{titlesec}
\usepackage{listings}
\usepackage[most]{tcolorbox}
\usepackage{url}
\urlstyle{tt}
\usepackage[hidelinks]{hyperref}
\setlength{\parskip}{0.5em}
\setlength{\parindent}{0pt}
\renewcommand{\arraystretch}{1.25}
\setlist{itemsep=0.2em,topsep=0.3em}
\definecolor{azul}{HTML}{1F3F8F}
\definecolor{rojo}{HTML}{C8102E}
\definecolor{azulclaro}{HTML}{3880FF}
\definecolor{codebg}{HTML}{F5F7FA}
\definecolor{codeframe}{HTML}{C8D0DC}
\pagestyle{fancy}\fancyhf{}
\fancyfoot[C]{\small P\'agina \thepage}
\fancyhead[L]{\small Pok\'edex con Ionic}
\renewcommand{\headrulewidth}{0.3pt}
\titleformat{\section}{\Large\bfseries\color{azul}}{\thesection}{0.6em}{}
\titleformat{\subsection}{\large\bfseries\color{azul}}{\thesubsection}{0.6em}{}
\lstdefinestyle{codigo}{
  basicstyle=\ttfamily\footnotesize,backgroundcolor=\color{codebg},frame=single,
  rulecolor=\color{codeframe},breaklines=true,columns=fullflexible,keepspaces=true,
  showstringspaces=false,upquote=true,xleftmargin=0.3em,xrightmargin=0.3em,aboveskip=0.4em,belowskip=0.6em,
  literate={á}{{\'a}}1 {é}{{\'e}}1 {í}{{\'i}}1 {ó}{{\'o}}1 {ú}{{\'u}}1 {ñ}{{\~n}}1 {Á}{{\'A}}1 {É}{{\'E}}1 {Í}{{\'I}}1 {Ó}{{\'O}}1 {Ú}{{\'U}}1 {¿}{{?`}}1 {¡}{{!`}}1}
\lstset{style=codigo}
\newcommand{\cmd}[1]{\texttt{#1}}
\newcommand{\archivo}[1]{\par\smallskip\noindent\textbf{Archivo:} \texttt{\detokenize{#1}}\par\nopagebreak}
\newtcolorbox{nota}[1][Nota]{breakable,colback=azulclaro!6,colframe=azulclaro!70!black,fonttitle=\bfseries,title=#1,boxrule=0.6pt,arc=2pt}
\newtcolorbox{alerta}[1][Importante]{breakable,colback=orange!8,colframe=orange!80!black,fonttitle=\bfseries,title=#1,boxrule=0.6pt,arc=2pt}
\newtcolorbox{resultado}[1][Resultado esperado]{breakable,colback=green!5,colframe=green!45!black,fonttitle=\bfseries,title=#1,boxrule=0.6pt,arc=2pt}

\title{{\Huge\bfseries\color{rojo} Pok\'edex con Ionic}\\[0.5em]{\Large Pr\'actica guiada en 2 etapas}}
\author{Andrick Millan y Emiliano Cisneros}
\date{\today}
\begin{document}\sloppy
\maketitle\thispagestyle{fancy}

\section{Qu\'e vamos a construir}
Una app de Pok\'emon con dos pantallas: una \textbf{lista} con im\'agenes y, al tocar una tarjeta,
una pantalla de \textbf{detalle} con tipos, habilidades, altura y peso que se piden a internet (PokeAPI).

\begin{center}
\begin{tabularx}{\linewidth}{@{}>{\bfseries}p{1.6cm} >{\raggedright\arraybackslash}X >{\raggedright\arraybackslash}p{4.4cm} r@{}}
\toprule
Etapa & Qu\'e se logra & Concepto nuevo & Min \\
\midrule
0 & Crear el proyecto y verlo en el navegador & \cmd{ionic start}, \cmd{ionic serve} & 10 \\
1 & Lista de 3 Pok\'emon con imagen & Componentes de Ionic, \cmd{@for} & 20 \\
2 & Pantalla de detalle con datos de internet & Rutas, \cmd{HttpClient}, signals & 30 \\
\bottomrule
\end{tabularx}
\end{center}

\begin{nota}[C\'omo usar esta gu\'ia]
En cada etapa: abre o crea los archivos indicados, \textbf{reemplaza todo su contenido} con el c\'odigo de la gu\'ia,
guarda con \texttt{Ctrl+S} y mira el navegador (se actualiza solo).
\end{nota}

\begin{alerta}[Reglas para evitar errores]
\begin{itemize}
\item Crea los archivos \textbf{a mano} (clic derecho $\rightarrow$ New File). No uses \cmd{ionic generate}.
\item Los nombres de archivo deben ser \textbf{exactamente} los de la gu\'ia (min\'usculas, con punto).
\item Crea el proyecto en una ruta corta, sin acentos ni espacios: \texttt{C:\textbackslash Proyectos} (Mac: \texttt{\~{}/Proyectos}).
\item Comentarios: en \texttt{.ts} se escribe \texttt{// texto}; en \texttt{.html} se escribe \texttt{<!-{}- texto -{}->}. Si pones \texttt{//} en un HTML, el texto aparece en la pantalla.
\item Las comillas especiales del c\'odigo (\texttt{`}) se llaman \emph{comillas invertidas}. Windows (teclado espa\~nol): tecla \texttt{`} y luego espacio, o \texttt{Alt Gr + \}}. Mac: \texttt{Option + `} y espacio.
\end{itemize}
\end{alerta}

\section{Etapa 0: preparar el proyecto (10 min)}
\subsection*{Requisitos previos}
\begin{itemize}
\item Node.js \textbf{22.22.3 o superior} (recomendado la LTS m\'as reciente). Comprobar: \cmd{node -v}
\item Ionic CLI: \cmd{npm install -g @ionic/cli}. Comprobar: \cmd{ionic -{}-version}
\item Visual Studio Code.
\end{itemize}

\subsection*{Pasos}
\begin{enumerate}
\item En VS Code abre una terminal (\texttt{Ctrl+`}) y entra a tu carpeta de trabajo:
\begin{lstlisting}
cd C:\Proyectos          # Mac: cd ~/Proyectos
\end{lstlisting}
\item Crea el proyecto (si pregunta por una cuenta de Ionic, responde \textbf{n}):
\begin{lstlisting}
ionic start pokedex blank --type=angular-standalone --capacitor
cd pokedex
\end{lstlisting}
\item Abre la carpeta del proyecto en VS Code: \texttt{File $\rightarrow$ Open Folder $\rightarrow$ pokedex}.
\item Arranca la app:
\begin{lstlisting}
ionic serve
\end{lstlisting}
\end{enumerate}
\begin{resultado}
Se abre \url{http://localhost:8100} con la p\'agina de bienvenida. \textbf{No cierres la terminal}: mientras
\cmd{ionic serve} corra, cada cambio guardado se ve solo. Para vista de celular: \texttt{F12} y el icono de tel\'efono.
\end{resultado}
\begin{nota}[Verifica esto]
Abre \texttt{src/app/home/home.page.ts}. Los componentes de Ionic deben importarse de \texttt{'@ionic/angular'}.
Si dice \texttt{'@ionic/angular/standalone'}, c\'ambialo por \texttt{'@ionic/angular'} en todos los c\'odigos de esta gu\'ia.
\end{nota}

\section{Etapa 1: lista de Pok\'emon (20 min)}
\textbf{Objetivo:} mostrar 3 Pok\'emon escritos a mano para entender c\'omo se arma una pantalla con Ionic.

\subsection*{Pasos}
Reemplaza el contenido de los tres archivos de \texttt{src/app/home/}.
\archivo{src/app/home/home.page.ts}
\begin{lstlisting}
import { Component } from '@angular/core';
import {
  IonHeader,
  IonToolbar,
  IonTitle,
  IonContent,
  IonCard,
  IonCardHeader,
  IonCardTitle,
  IonCardSubtitle,
} from '@ionic/angular';

@Component({
  selector: 'app-home',
  templateUrl: 'home.page.html',
  styleUrls: ['home.page.scss'],
  imports: [
    IonHeader,
    IonToolbar,
    IonTitle,
    IonContent,
    IonCard,
    IonCardHeader,
    IonCardTitle,
    IonCardSubtitle,
  ],
})
export class HomePage {
  // Lista escrita a mano (todavia no usamos internet)
  pokemons = [
    { id: 1, nombre: 'bulbasaur' },
    { id: 4, nombre: 'charmander' },
    { id: 25, nombre: 'pikachu' },
  ];

  // Devuelve la URL de la imagen de un Pokemon segun su numero
  imagen(id: number): string {
    return `https://raw.githubusercontent.com/PokeAPI/sprites/master/sprites/pokemon/other/official-artwork/${id}.png`;
  }
}
\end{lstlisting}
\archivo{src/app/home/home.page.html}
\begin{lstlisting}
<ion-header>
  <ion-toolbar color="danger">
    <ion-title>Pokédex</ion-title>
  </ion-toolbar>
</ion-header>

<ion-content class="ion-padding">
  @for (p of pokemons; track p.id) {
    <ion-card>
      <img [src]="imagen(p.id)" [alt]="p.nombre" />
      <ion-card-header>
        <ion-card-subtitle>#{{ p.id }}</ion-card-subtitle>
        <ion-card-title>{{ p.nombre }}</ion-card-title>
      </ion-card-header>
    </ion-card>
  }
</ion-content>
\end{lstlisting}
\archivo{src/app/home/home.page.scss}
\begin{lstlisting}
/* Las imagenes ocupan todo el ancho de la tarjeta */
img {
  width: 100%;
  height: auto;
  display: block;
}
\end{lstlisting}

\begin{resultado}
Tres tarjetas con la imagen de Bulbasaur, Charmander y Pikachu y una barra roja con el t\'itulo \emph{Pok\'edex}.
\end{resultado}
\subsection*{Qu\'e aprendimos}
\begin{itemize}
\item Una \textbf{p\'agina} son tres archivos: \texttt{.ts} (l\'ogica), \texttt{.html} (lo que se ve) y \texttt{.scss} (estilo).
\item \texttt{ion-header}, \texttt{ion-toolbar}, \texttt{ion-content}, \texttt{ion-card}: componentes ya hechos de Ionic. Cada uno que se usa en el HTML debe estar en la lista \texttt{imports} del componente \textbf{y} en el \texttt{import} de arriba.
\item \cmd{@for} repite un bloque por cada elemento de la lista; \cmd{track} indica qu\'e dato identifica a cada elemento.
\item \cmd{[src]="..."} conecta un valor del c\'odigo con el HTML; \cmd{\{\{ ... \}\}} imprime un valor.
\item Cada n\'umero es el de la Pok\'edex real: 1 Bulbasaur, 4 Charmander, 25 Pikachu.
\end{itemize}
\begin{alerta}[Errores comunes]
\begin{center}
\begin{tabularx}{\linewidth}{@{}>{\raggedright\arraybackslash}p{5cm} >{\raggedright\arraybackslash}X@{}}
Un componente sale en gris en la lista de imports & Est\'a importado arriba pero falta en \texttt{imports: [...]}. \\
\texttt{'ion-card' is not a known element} & Agrega \texttt{IonCard} a \texttt{imports: [...]}. \\
Im\'agenes rotas & La URL debe ir entre \textbf{comillas invertidas}, sin espacios, y terminar en \texttt{\$\{id\}.png}. \\
Pantalla en blanco & Revisa la terminal y la consola (\texttt{F12}); suele ser una llave sin cerrar. \\
\end{tabularx}
\end{center}
\end{alerta}

\section{Etapa 2: pantalla de detalle (30 min)}
\textbf{Objetivo:} al tocar una tarjeta, abrir otra pantalla con los datos de ese Pok\'emon, pedidos a la API.

\subsection*{Idea}
Una \textbf{API} es un servicio al que le pedimos datos por internet. Abre \url{https://pokeapi.co/api/v2/pokemon/25}
en el navegador: es un JSON con todo lo de Pikachu (\texttt{types}, \texttt{abilities}, \texttt{height}, \texttt{weight}).
Nuestra pantalla de detalle pide ese JSON y muestra las partes que nos interesan.

\subsection*{Pasos}
\begin{enumerate}
\item \textbf{Permitir peticiones a internet.} En \texttt{src/main.ts} agrega el \texttt{import} y la l\'inea \cmd{provideHttpClient()} dentro de \texttt{providers}. Puedes reemplazar todo el archivo:
\end{enumerate}
\archivo{src/main.ts}
\begin{lstlisting}
import { bootstrapApplication } from '@angular/platform-browser';
import { RouteReuseStrategy, provideRouter, withComponentInputBinding, withPreloading, PreloadAllModules } from '@angular/router';
import { provideHttpClient } from '@angular/common/http';
import { IonicRouteStrategy, provideIonicAngular } from '@ionic/angular';

import { routes } from './app/app.routes';
import { AppComponent } from './app/app.component';

bootstrapApplication(AppComponent, {
  providers: [
    { provide: RouteReuseStrategy, useClass: IonicRouteStrategy },
    provideIonicAngular(),
    provideRouter(routes, withPreloading(PreloadAllModules), withComponentInputBinding()),
    provideHttpClient(), // <- permite consultar la PokeAPI
  ],
});
\end{lstlisting}

\begin{enumerate}\setcounter{enumi}{1}
\item \textbf{Agregar la ruta} de la pantalla nueva (\texttt{detalle/:id}; los dos puntos indican que el n\'umero cambia). Reemplaza \texttt{src/app/app.routes.ts}:
\end{enumerate}
\archivo{src/app/app.routes.ts}
\begin{lstlisting}
import { Routes } from '@angular/router';

export const routes: Routes = [
  {
    path: 'home',
    loadComponent: () => import('./home/home.page').then((m) => m.HomePage),
  },
  {
    path: 'detalle/:id',
    loadComponent: () => import('./detalle/detalle.page').then((m) => m.DetallePage),
  },
  {
    path: '',
    redirectTo: 'home',
    pathMatch: 'full',
  },
];
\end{lstlisting}

\begin{enumerate}\setcounter{enumi}{2}
\item \textbf{Hacer que la tarjeta sea un enlace.} En \texttt{home.page.ts} se agrega \texttt{RouterLink} (en el \texttt{import} de \texttt{@angular/router} \textbf{y} dentro de \texttt{imports: [...]}; no va en el bloque de Ionic). En \texttt{home.page.html} la tarjeta lleva \cmd{[routerLink]}. Reemplaza ambos archivos:
\end{enumerate}
\archivo{src/app/home/home.page.ts}
\begin{lstlisting}
import { Component } from '@angular/core';
import { RouterLink } from '@angular/router';
import {
  IonHeader,
  IonToolbar,
  IonTitle,
  IonContent,
  IonCard,
  IonCardHeader,
  IonCardTitle,
  IonCardSubtitle,
} from '@ionic/angular';

@Component({
  selector: 'app-home',
  templateUrl: 'home.page.html',
  styleUrls: ['home.page.scss'],
  imports: [
    RouterLink,
    IonHeader,
    IonToolbar,
    IonTitle,
    IonContent,
    IonCard,
    IonCardHeader,
    IonCardTitle,
    IonCardSubtitle,
  ],
})
export class HomePage {
  pokemons = [
    { id: 1, nombre: 'bulbasaur' },
    { id: 4, nombre: 'charmander' },
    { id: 25, nombre: 'pikachu' },
  ];

  imagen(id: number): string {
    return `https://raw.githubusercontent.com/PokeAPI/sprites/master/sprites/pokemon/other/official-artwork/${id}.png`;
  }
}
\end{lstlisting}
\archivo{src/app/home/home.page.html}
\begin{lstlisting}
<ion-header>
  <ion-toolbar color="danger">
    <ion-title>Pokédex</ion-title>
  </ion-toolbar>
</ion-header>

<ion-content class="ion-padding">
  @for (p of pokemons; track p.id) {
    <ion-card [routerLink]="['/detalle', p.id]">
      <img [src]="imagen(p.id)" [alt]="p.nombre" />
      <ion-card-header>
        <ion-card-subtitle>#{{ p.id }}</ion-card-subtitle>
        <ion-card-title>{{ p.nombre }}</ion-card-title>
      </ion-card-header>
    </ion-card>
  }
</ion-content>
\end{lstlisting}

\begin{enumerate}\setcounter{enumi}{3}
\item \textbf{Crear la pantalla de detalle.} Crea la carpeta \texttt{src/app/detalle/} (clic derecho en \texttt{app} $\rightarrow$ New Folder) y dentro dos archivos nuevos:
\end{enumerate}
\archivo{src/app/detalle/detalle.page.ts}
\begin{lstlisting}
import { Component, inject, signal } from '@angular/core';
import { ActivatedRoute } from '@angular/router';
import { HttpClient } from '@angular/common/http';
import {
  IonHeader,
  IonToolbar,
  IonButtons,
  IonBackButton,
  IonTitle,
  IonContent,
  IonChip,
  IonSpinner,
} from '@ionic/angular';

@Component({
  selector: 'app-detalle',
  templateUrl: 'detalle.page.html',
  imports: [
    IonHeader,
    IonToolbar,
    IonButtons,
    IonBackButton,
    IonTitle,
    IonContent,
    IonChip,
    IonSpinner,
  ],
})
export class DetallePage {
  private ruta = inject(ActivatedRoute);
  private http = inject(HttpClient);

  // Aquí guardamos los datos que llegan de internet (al inicio no hay nada)
  pokemon = signal<any>(null);

  constructor() {
    // Leemos el número que viene en la dirección: /detalle/25
    const id = this.ruta.snapshot.paramMap.get('id');

    // Pedimos los datos de ese Pokémon a la PokéAPI
    this.http
      .get<any>(`https://pokeapi.co/api/v2/pokemon/${id}`)
      .subscribe((datos) => this.pokemon.set(datos));
  }

  imagen(id: number): string {
    return `https://raw.githubusercontent.com/PokeAPI/sprites/master/sprites/pokemon/other/official-artwork/${id}.png`;
  }
}
\end{lstlisting}
\archivo{src/app/detalle/detalle.page.html}
\begin{lstlisting}
<ion-header>
  <ion-toolbar color="danger">
    <ion-buttons slot="start">
      <ion-back-button defaultHref="/home"></ion-back-button>
    </ion-buttons>
    <ion-title>Detalle</ion-title>
  </ion-toolbar>
</ion-header>

<ion-content class="ion-padding">
  @if (pokemon(); as p) {
    <div class="ion-text-center">
      <img [src]="imagen(p.id)" [alt]="p.name" style="width: 220px; max-width: 70%" />
      <h1>{{ p.name }}</h1>
    </div>

    <h3>Tipos</h3>
    @for (t of p.types; track t.type.name) {
      <ion-chip>{{ t.type.name }}</ion-chip>
    }

    <h3>Habilidades</h3>
    @for (h of p.abilities; track h.ability.name) {
      <ion-chip color="danger">{{ h.ability.name }}</ion-chip>
    }

    <h3>Medidas</h3>
    <p>Altura: {{ p.height / 10 }} m</p>
    <p>Peso: {{ p.weight / 10 }} kg</p>
  } @else {
    <div class="ion-text-center">
      <ion-spinner></ion-spinner>
    </div>
  }
</ion-content>
\end{lstlisting}

\begin{enumerate}\setcounter{enumi}{4}
\item Guarda todo (\texttt{File $\rightarrow$ Save All}). Si la terminal muestra un error viejo, deten el servidor con \texttt{Ctrl+C} y ejecuta de nuevo \cmd{ionic serve}.
\end{enumerate}

\begin{resultado}
Al tocar la tarjeta de Pikachu, la URL cambia a \texttt{/detalle/25} y se ve el Pok\'emon grande con sus tipos (\emph{electric}),
sus habilidades (\emph{static}, \emph{lightning-rod}), altura 0.4~m y peso 6~kg. La flecha de regreso vuelve a la lista.
\end{resultado}
\subsection*{Qu\'e aprendimos}
\begin{itemize}
\item \textbf{Ruta con par\'ametro:} \texttt{detalle/:id}; el n\'umero se lee con \cmd{ActivatedRoute}.
\item \cmd{routerLink}: navegar sin recargar la p\'agina.
\item \cmd{HttpClient}: pedir datos a una API.
\item \textbf{signal:} variable ``reactiva''. Se lee con par\'entesis, \cmd{pokemon()}, y se cambia con \cmd{.set(...)}; cuando cambia, la pantalla se actualiza sola.
\item \cmd{@if / @else}: mostrar cosas seg\'un una condici\'on (aqu\'i, el \emph{spinner} mientras llegan los datos).
\end{itemize}
\begin{alerta}[Errores comunes]
\begin{center}
\begin{tabularx}{\linewidth}{@{}>{\raggedright\arraybackslash}p{5.3cm} >{\raggedright\arraybackslash}X@{}}
\texttt{Could not resolve ./detalle/detalle.page} & El archivo \texttt{detalle.page.ts} no existe, est\'a vac\'io o con otro nombre. Si ya est\'a bien, reinicia \cmd{ionic serve}. \\
\texttt{Can't bind to 'routerLink'} & Falta \texttt{RouterLink} en \texttt{imports: [...]} de \texttt{home.page.ts}. \\
\texttt{No provider for HttpClient} & Falta \cmd{provideHttpClient()} en \texttt{main.ts}. \\
Sale texto raro con \texttt{//} en la pantalla & Un comentario de HTML debe ser \texttt{<!-{}- ... -{}->}. \\
Spinner infinito & Sin internet o URL mal escrita. Revisa la consola (\texttt{F12}). \\
Ves \texttt{RouterLink} en rojo & Est\'a dentro del bloque \texttt{from '@ionic/angular'}; debe ir en su propio \texttt{import} de \texttt{'@angular/router'}. \\
\end{tabularx}
\end{center}
\end{alerta}

\section{Cierre y retos}
\begin{itemize}
\item \textbf{Reto 1:} agrega otro Pok\'emon a la lista (por ejemplo, 7 Squirtle o 150 Mewtwo).
\item \textbf{Reto 2:} en el detalle, muestra tambi\'en el \texttt{base\_experience} del Pok\'emon (est\'a en el JSON).
\item \textbf{Reto 3:} cambia el color de la barra superior (\texttt{color="danger"}) por \texttt{primary}, \texttt{success} o \texttt{warning}.
\item (Opcional, Android) \cmd{ionic build}, \cmd{npx cap add android}, \cmd{npx cap open android}: la misma app corre como app nativa.
\end{itemize}

\section{Soluci\'on de problemas generales}
\begin{center}
\begin{tabularx}{\linewidth}{@{}>{\raggedright\arraybackslash}p{5.2cm} >{\raggedright\arraybackslash}X@{}}
\toprule\textbf{Problema} & \textbf{Soluci\'on}\\\midrule
La terminal no reconoce \cmd{ionic} & Cierra y vuelve a abrir VS Code; si sigue, reinstala con \cmd{npm install -g @ionic/cli}. \\
PowerShell bloquea scripts & \cmd{Set-ExecutionPolicy -Scope CurrentUser RemoteSigned} \\
Errores raros de compilaci\'on en Windows & Crea el proyecto en \texttt{C:\textbackslash Proyectos}, sin acentos ni carpetas redirigidas. \\
Puerto 8100 ocupado & \cmd{ionic serve -{}-port=8101} \\
Guard\'e y no cambia nada & Revisa que \cmd{ionic serve} siga corriendo; si hay un error viejo, \texttt{Ctrl+C} y vuelve a ejecutarlo. \\
Formatear el c\'odigo en VS Code & \texttt{Shift+Alt+F} (Mac: \texttt{Shift+Option+F}). \\
\bottomrule\end{tabularx}
\end{center}
\end{document}